\documentclass[11pt,a4paper]{article}
\usepackage[utf8]{inputenc}
\usepackage[T1]{fontenc}
\usepackage{amsmath,amssymb}
\usepackage{hyperref}
\usepackage{listings}
\usepackage{xcolor}
\usepackage{geometry}
\geometry{margin=2.5cm}

\newcommand\YAMLkeystyle{\color{black}\bfseries}
\newcommand\YAMLvaluestyle{\color{blue}\mdseries}

\lstdefinelanguage{yaml}
{
  keywords={true,false,null,y,n},
  keywordstyle=\color{darkgray}\bfseries,
  basicstyle=\YAMLkeystyle,
  sensitive=false,
  comment=[l]{\#},
  morecomment=[s]{/*}{*/},
  commentstyle=\color{purple}\ttfamily,
  stringstyle=\YAMLvaluestyle\ttfamily,
  moredelim=[l][\color{orange}]{\&},
  moredelim=[l][\color{magenta}]{*},
  moredelim=**[il][\color{purple}]{:},
  morestring=[b]',
  morestring=[b]"
}

\lstset{
  basicstyle=\ttfamily\small,
  backgroundcolor=\color{gray!10},
  frame=single,
  breaklines=true,
  keywordstyle=\color{blue},
  commentstyle=\color{green!50!black},
  stringstyle=\color{red!70!black},
}

\title{Bayesian Decision Tree with Monte Carlo Simulation\\[0.5em]
\large Theory, Implementation, and YAML Configuration Guide}
\author{Ari-Pekka Sihvonen\\
\texttt{ap.sihvonen@gmail.com}}
\date{\copyright\ 2026 — MIT License}

\begin{document}
\maketitle

\begin{abstract}
This document describes the theory, implementation, and usage of
\texttt{bayes-tree-eng.py}, a tool that encodes hypotheses as Bayesian decision
trees and uses Monte Carlo simulation to propagate uncertainty. Top-level
branches act as independent evidence combined at the root via likelihood ratios
in log-odds space. The tool produces posterior distributions, sensitivity
analyses, and importance rankings.
\end{abstract}

\tableofcontents
\newpage

%═══════════════════════════════════════════════════════════════════════════════
\section{Introduction}

The Bayesian Decision Tree tool combines Bayesian inference with Monte Carlo
simulation to propagate uncertainty through a tree of evidence. Each branch
of the tree represents an independent piece of evidence characterized by a
\emph{likelihood ratio} (LR), and the tool computes a posterior probability
distribution for the root hypothesis.

The motivation is to provide a structured, transparent framework for combining
multiple heuristic likelihood ratios into a posterior probability. It is useful
for exploratory ``what-if'' analysis, sensitivity studies, and producing an
interpretable tree of how local evidence shifts beliefs.

%═══════════════════════════════════════════════════════════════════════════════
\section{Theoretical Background}

\subsection{Bayes' Theorem}

For a hypothesis $H$ and evidence $E$:
\begin{equation}
P(H|E) = \frac{P(E|H)\,P(H)}{P(E)}
\end{equation}

The \textbf{likelihood ratio} (LR) form is:
\begin{equation}
\text{LR} = \frac{P(E|H)}{P(E|\neg H)}
\end{equation}

\subsection{Log-Odds Representation}

Converting probability to log-odds:
\begin{equation}
\text{lo}(p) = \ln\frac{p}{1-p}
\end{equation}

And back:
\begin{equation}
p = \frac{1}{1 + e^{-\text{lo}}}
\end{equation}

A Bayesian update in log-odds space is simply additive:
\begin{equation}
\text{lo}(P(H|E)) = \text{lo}(P(H)) + \ln(\text{LR})
\end{equation}

This additive property makes it natural to combine multiple independent pieces
of evidence.

\subsection{Combining Independent Evidence}

Given $n$ independent evidence branches with likelihood ratios
$\text{LR}_1, \text{LR}_2, \ldots, \text{LR}_n$, the combined posterior
in log-odds is:
\begin{equation}
\text{lo}(\text{posterior}) = \text{lo}(\text{prior}) + \sum_{i=1}^{n} \ln(\text{LR}_i)
\end{equation}

Equivalently, in odds form:
\begin{equation}
\text{odds}(H \mid E_{1:n}) = \text{odds}(H) \prod_{i=1}^{n} \text{LR}_i
\end{equation}

The tool computes this at the root level, treating each top-level child as an
independent evidence source.

\subsection{Effective Likelihood Ratio}

The \emph{effective LR} is the single likelihood ratio that would produce the
same posterior from the given prior:
\begin{equation}
\text{LR}_{\text{eff}} = \exp\!\Big(\text{lo}(\text{posterior}) - \text{lo}(\text{prior})\Big)
\end{equation}

\subsection{Monte Carlo Uncertainty Propagation}

When the LR for each branch is uncertain (specified as an interval rather than
a point value), the tool uses Monte Carlo simulation:

\begin{enumerate}
\item For each simulation run, sample an LR from each branch's distribution.
\item Compute the combined posterior using the log-odds sum.
\item Repeat $N$ times (default: 10{,}000) to build a posterior distribution.
\item Report summary statistics: median, mean, 90\% credible interval.
\end{enumerate}

\subsection{LR Sampling Distributions}

Three distributions are available for sampling LR values from the interval
$[\text{lr\_min}, \text{lr\_max}]$:

\begin{description}
\item[log\_uniform] (default) --- Uniform in log-space. Best when uncertainty
  spans orders of magnitude. The sampled value is:
  \begin{equation}
  \text{LR} = \exp\!\Big(\text{Uniform}\big(\ln(\text{lr\_min}),\, \ln(\text{lr\_max})\big)\Big)
  \end{equation}

\item[uniform] --- Uniform in linear space:
  \begin{equation}
  \text{LR} = \text{Uniform}(\text{lr\_min},\, \text{lr\_max})
  \end{equation}

\item[beta] --- Beta distribution scaled to the interval, controlled by
  parameters \texttt{lr\_alpha} and \texttt{lr\_beta}:
  \begin{equation}
  \text{LR} = \text{lr\_min} + \text{Beta}(\alpha, \beta) \times (\text{lr\_max} - \text{lr\_min})
  \end{equation}
\end{description}

%═══════════════════════════════════════════════════════════════════════════════
\section{Tree Structure and Chaining}

The tree has two distinct roles:
\begin{itemize}
\item \textbf{Root level:} Children are combined as independent evidence
  (log-odds sum). Each child contributes its own LR to the root posterior.
\item \textbf{Sub-levels:} Children represent drill-down or decomposition of
  the parent evidence. The child's prior is set to the parent's posterior
  (chaining), showing how the probability evolves through successive reasoning
  steps.
\end{itemize}

Importantly, children of top-level branches are \emph{not} multiplied into the
root combination. They only serve to explain or decompose a branch's LR. If
child-level evidence should aggregate to the root, the model must be
restructured so that the relevant nodes appear at the top level.

%═══════════════════════════════════════════════════════════════════════════════
\section{Sensitivity Analysis}

The tool provides two sensitivity measures:
\begin{enumerate}
\item \textbf{Threshold probabilities:} $P(\text{posterior} < 5\%)$,
  $P(\text{posterior} < 10\%)$, $P(\text{posterior} > 50\%)$.
\item \textbf{Leave-one-out ranking:} For each branch, compute the posterior
  without that branch. The difference from the baseline median shows how much
  each branch influences the result.
\end{enumerate}

%═══════════════════════════════════════════════════════════════════════════════
\section{YAML File Structure}

The input is a YAML file defining a tree of evidence. Below is the complete
specification.

\subsection{Root Node}

\begin{lstlisting}[language=yaml]
node: "Hypothesis name"
prior: 0.5                  # Prior probability (0 to 1)
children:
  - ...                     # List of evidence branches
\end{lstlisting}

\subsection{Evidence Node}

Each child node represents one piece of evidence:

\begin{lstlisting}[language=yaml]
- node: "Evidence description"
  evidence_type: for        # "for", "against", or "neutral"

  # Option A: Point estimate
  likelihood_ratio: 2.5     # Exact LR value

  # Option B: Uncertainty interval (preferred)
  lr_min: 0.01              # Lower bound of LR
  lr_max: 0.10              # Upper bound of LR
  lr_dist: log_uniform      # Distribution: log_uniform | uniform | beta

  # Optional for beta distribution
  lr_alpha: 2               # Beta shape parameter alpha
  lr_beta: 2                # Beta shape parameter beta

  # Optional sub-evidence (drill-down)
  children:
    - node: "Sub-evidence"
      evidence_type: for
      likelihood_ratio: 1.5
\end{lstlisting}

\subsection{Evidence Type}

\begin{description}
\item[for] --- Evidence supporting the hypothesis ($\text{LR} > 1$).
\item[against] --- Evidence against the hypothesis ($\text{LR} < 1$).
\item[neutral] --- Uninformative evidence ($\text{LR} \approx 1$).
\end{description}

The tool validates that \texttt{evidence\_type} is consistent with the LR
values and warns if there is a mismatch.

\subsection{Likelihood Ratio Interpretation}

\begin{center}
\begin{tabular}{ll}
\hline
\textbf{LR value} & \textbf{Meaning} \\
\hline
LR $\gg$ 1 & Strong evidence for the hypothesis \\
LR $>$ 1   & Moderate evidence for \\
LR $=$ 1   & Neutral (no update) \\
LR $<$ 1   & Evidence against \\
LR $\ll$ 1 & Strong evidence against \\
\hline
\end{tabular}
\end{center}

%═══════════════════════════════════════════════════════════════════════════════
\section{Complete YAML Example}

\begin{lstlisting}[language=yaml]
node: "Patient has disease X"
prior: 0.10
children:
  - node: "Positive blood test"
    evidence_type: for
    lr_min: 3.0
    lr_max: 8.0
    lr_dist: log_uniform
    children:
      - node: "Test specificity concerns"
        evidence_type: against
        lr_min: 0.5
        lr_max: 0.9

  - node: "Family history present"
    evidence_type: for
    likelihood_ratio: 2.0

  - node: "Age below typical onset"
    evidence_type: against
    lr_min: 0.3
    lr_max: 0.7
    lr_dist: uniform

  - node: "Symptom pattern matches"
    evidence_type: for
    lr_min: 1.5
    lr_max: 4.0
    lr_dist: beta
    lr_alpha: 2
    lr_beta: 3
\end{lstlisting}

%═══════════════════════════════════════════════════════════════════════════════
\section{Usage}

\subsection{Running the Tool}

\begin{lstlisting}[language=bash]
# Default: 10,000 simulations
python bayes-tree-eng.py examples/mycase.yaml

# Custom number of simulations
python bayes-tree-eng.py examples/mycase.yaml 50000
\end{lstlisting}

\subsection{Output Sections}

The tool produces the following output:
\begin{enumerate}
\item \textbf{Validation warnings} --- Flags any inconsistencies between
  \texttt{evidence\_type} and LR values.
\item \textbf{Combined posterior} --- Summary statistics (median, mean, CI)
  and ASCII histogram of the posterior distribution.
\item \textbf{Tree view} --- Hierarchical display showing each branch's
  contribution, with prior $\to$ posterior transitions.
\item \textbf{Sensitivity analysis} --- Threshold probabilities.
\item \textbf{Importance ranking} --- Leave-one-out analysis showing which
  branches have the most influence.
\end{enumerate}

\subsection{Tips for Creating YAML Files}

\begin{enumerate}
\item \textbf{Start with the prior.} What is the base rate or initial
  probability before any evidence is considered?
\item \textbf{List independent evidence.} Each top-level child should be
  logically independent of the others.
\item \textbf{Prefer intervals over point estimates.} Using \texttt{lr\_min}
  and \texttt{lr\_max} captures your uncertainty about the strength of
  evidence.
\item \textbf{Use log\_uniform for uncertain LRs.} When you know the LR is
  ``somewhere between 0.01 and 0.1'' but aren't sure where, log-uniform
  gives appropriate weight to different orders of magnitude.
\item \textbf{Add children for drill-down.} Sub-nodes explain or decompose
  a parent's evidence but don't add to the root combination.
\item \textbf{Validate your intuitions.} Run the tool and check that the
  output matches your qualitative expectation. If not, reconsider your LR
  estimates.
\end{enumerate}

%═══════════════════════════════════════════════════════════════════════════════
\section{Limitations and Assumptions}

\begin{itemize}
\item \textbf{Independence assumption.} Top-level branches are assumed
  conditionally independent given $H$ and $\neg H$. This is a strong
  assumption and should be scrutinized when modeling correlated evidence.
  In practice, if two branches share a common information source, they
  should be merged or one should be made a sub-node of the other.

\item \textbf{Children do not aggregate to root.} Children are used only for
  chained drill-down display. If child-level evidence should contribute to
  the root posterior, the model must be restructured so that the relevant
  nodes appear at the top level.

\item \textbf{Log-uniform bias.} The default log-uniform sampling biases
  toward multiplicative scales. Users should choose a distribution that
  reflects their actual beliefs about where in the interval the true LR lies.

\item \textbf{Subjective LR estimates.} The tool is only as good as the
  likelihood ratios you assign. It makes subjective reasoning transparent
  and quantitative, but does not eliminate subjectivity.

\item \textbf{Numerical stability.} The implementation clamps probabilities
  and log-odds to safe numerical ranges to avoid overflow, which may
  slightly affect extreme posteriors (very close to 0 or 1).
\end{itemize}

%═══════════════════════════════════════════════════════════════════════════════
\section{Implementation Notes}

\begin{itemize}
\item \textbf{Log-odds stability:} Probabilities are clamped to
  $[10^{-12},\, 1-10^{-12}]$ before log-odds conversion; log-odds are
  clamped to $[-700, 700]$ before exponentiation.
\item \textbf{Sampling:} The same random draw strategy is applied per
  simulation; users can set simulation counts via the command line.
\item \textbf{Tree view:} A recursive collection routine computes node-level
  posteriors by feeding a parent's posterior median as prior into children,
  producing local medians and credible intervals.
\item \textbf{Effective LR for root:} The root node's LR interval is derived
  from the 5th--95th percentiles of the effective LR distribution across all
  Monte Carlo runs.
\end{itemize}

%═══════════════════════════════════════════════════════════════════════════════
\section{Dependencies}

The Python script requires:
\begin{itemize}
\item Python 3.8+
\item \texttt{PyYAML} (\texttt{pip install pyyaml})
\end{itemize}

\end{document}
